\chapter{Conclusion}
\label{ch:sum}

This thesis presented the design and implementation of an
\emph{AI Cookbook with Ingredient-Based Search}, a web application that
combines personalized recipe generation, OCR-based nutrition-label
processing, and persistent user-specific history. The project was
motivated by two practical problems: deciding what to cook from
available ingredients under dietary constraints, and converting
nutrition-label information from food packaging into structured digital
data without manual transcription.

The final system addresses these problems through a full-stack web
application architecture. Users can submit ingredient lists and receive
structured recipes grounded by retrieval-augmented generation, while
food-label images can be processed through an OCR pipeline that returns
both raw extracted text and structured nutritional fields.
Authenticated users can additionally store dietary defaults, save
generated recipes, and preserve OCR extraction history for later review.

\section{Summary of achieved goals versus original proposal commitments}

The project achieved the main objectives defined in the introduction and
covered the core commitments of the original proposal.

\begin{itemize}
	\item A working recipe-generation workflow was implemented with support
	for ingredients, dietary preferences, allergy constraints, retrieval
	context, warnings, and structured output.
	\item A user account system was implemented with registration, login,
	profile-default storage, recipe history, and saved-recipe support.
	\item An OCR-based nutrition-label workflow was implemented with image
	validation, preprocessing, OCR extraction, structured nutrition
	conversion, and optional persistence.
	\item The system was supported by a real persistence layer, migration
	history, vector retrieval subsystem, and media handling logic.
	\item A multi-layer testing strategy was prepared combining automated
	tests, manual workflow tests, and AI-specific evaluation materials.
\end{itemize}

Taken together, these results show that the final system is a working
application with multiple integrated features rather than an isolated
demonstration of a language-model call.

\section{Lessons learned and challenges encountered}

During the project, the main technical challenge was integrating several
different concerns into one coherent application:

\begin{itemize}
	\item retrieval-augmented generation for recipe personalization,
	\item rule-supported safety filtering for allergies and dietary
	preferences,
	\item OCR combined with language-model-based semantic structuring,
	\item persistence of both classical user data and AI-generated
	content,
	\item a modular architecture separating routes, services, models,
	schemas, and provider abstractions.
\end{itemize}

One of the main lessons learned was that AI-assisted features required
supporting software-engineering structure around them. Validation,
persistence, testing, and explicit limitation handling were necessary
for turning recipe generation and OCR extraction into stable application
features rather than isolated API calls.

Another lesson was that system boundaries had to remain explicit.
External AI providers made output quality and availability partly
dependent on services outside the local application. This affected both
implementation and testing choices, especially for OCR extraction,
recipe generation, and evaluation of output quality.

\section{Current limitations}

Despite the achieved results, the final system still has several
important limitations.

\begin{itemize}
	\item The quality of recipe generation, OCR text extraction, and
	structured nutrition output depends on external AI providers.
	\item The nutritional estimates attached to generated recipes are
	approximate and should not be interpreted as certified dietary
	analysis.
	\item Allergy and dietary handling is safety-oriented but not
	medically complete, since it is based on rule-based filtering and
	prompt-level constraints rather than formal nutritional reasoning.
	\item The default storage and media-serving model is appropriate for
	local thesis deployment, but it is not intended as a final production
	deployment architecture.
	\item The current vector retrieval corpus is intentionally limited in
	scope and therefore cannot represent every culinary style or dietary
	context.
\end{itemize}

These limitations define the boundary between what the system currently
demonstrates and what would still be required for broader real-world
deployment.

\section{Proposed future work and development directions}

The following improvements would be natural continuations of the
project:

\begin{itemize}
	\item Expanding the recipe corpus and improving retrieval quality with
	better metadata and reranking strategies.
	\item Introducing richer user modelling, such as cuisine preferences,
	cooking-time preferences, or meal-type preferences.
	\item Replacing local media storage and process-local rate limiting
	with more deployment-ready infrastructure.
	\item Improving OCR robustness through broader preprocessing options
	and more diverse evaluation datasets.
	\item Extending the nutritional model beyond a small fixed schema
	toward richer product comparison and dietary tracking features.
\end{itemize}

These directions would extend the system beyond the current thesis
scope while building directly on the implemented architecture and
testing results.
